% Load the beamer-sharcnet document class...
\documentclass[tabu]{beamer-uwindsor}

% Have latex inform when another run is needed...
\usepackage{rerunfilecheck}

% Use the csquotes package...
\usepackage{csquotes}

% Use the listings package...
\usepackage{listings}

\lstloadlanguages{Haskell}

\lstnewenvironment{code}
{\lstset{}%
	\csname lst@SetFirstLabel\endcsname}
{\csname lst@SaveFirstLabel\endcsname}
\lstset{
	basicstyle={\scriptsize\ttfamily\singlespacing},
	flexiblecolumns=false,
	basewidth={0.5em,0.45em},
	literate={+}{{$+$}}1 {/}{{$/$}}1 {*}{{$*$}}1 {=}{{$=$}}1
	{>}{{$>$}}1 {<}{{$<$}}1 {\\}{{$\lambda$}}1
	{\\\\}{{\char`\\\char`\\}}1
	{->}{{$\rightarrow$}}2 {>=}{{$\geq$}}2 {<-}{{$\leftarrow$}}2
	{<=}{{$\leq$}}2 {=>}{{$\Rightarrow$}}2
	{\ .\ }{{$\circ$}}2
	{>>}{{>>}}2 {>>=}{{>>=}}2
	{|}{{$\mid$}}1
}

\lstnewenvironment{spec}
{\lstset{}%
	\csname lst@SetFirstLabel\endcsname}
{\csname lst@SaveFirstLabel\endcsname}
\lstset{
	basicstyle={\scriptsize\ttfamily\singlespacing},
	flexiblecolumns=false,
	basewidth={0.5em,0.45em},
	literate={+}{{$+$}}1 {/}{{$/$}}1 {*}{{$*$}}1 {=}{{$=$}}1
	{>}{{$>$}}1 {<}{{$<$}}1 {\\}{{$\lambda$}}1
	{\\\\}{{\char`\\\char`\\}}1
	{->}{{$\rightarrow$}}2 {>=}{{$\geq$}}2 {<-}{{$\leftarrow$}}2
	{<=}{{$\leq$}}2 {=>}{{$\Rightarrow$}}2
	{\ .\ }{{$\circ$}}2
	{>>}{{>>}}2 {>>=}{{>>=}}2
	{|}{{$\mid$}}1
}

%
% Use TrueType fonts (i.e., xelatex).
%
% Use "Charis SIL" font even though it is not sans-serif as
% it is easy-to-read, has all Unicode characters, is under
% the open SIL font licence, and is free.
%
% As SIL does not have a monospaced font, Liberation Mono
% is used instead.
%
\usepackage{xltxtra,fontspec,xunicode}
\defaultfontfeatures{Scale=MatchLowercase}
\setromanfont{Charis SIL}
\setsansfont{Charis SIL}
\setmainfont{Charis SIL}
%\setmonofont[Scale=0.70]{Liberation Mono}
\setmonofont{Liberation Mono}

%
% Define slide show and presenter-specific globals...
%
\providecommand{\basetexfilename}{slides}
\providecommand{\emailshane}{peelar@uwindsor.ca}
\title[Masters Thesis Defence]{Accommodating prepositional phrases in a highly modular
	natural language query interface to semantic web triplestores using a novel event-based denotational semantics for English and a set of parser combinators.}
\date{December 12th 2016}
\author[Shane Peelar]{%
  Shane Peelar, BSc (Hons) \texorpdfstring{\\}{}%
  {\footnotesize\href{mailto:\emailshane}{\emailshane}} \texorpdfstring{\\}{}%
}
\institute[University of Windsor]{%
  University of Windsor, \texorpdfstring{\\}{}%
  Windsor, Ontario, Canada \texorpdfstring{\\}{}%
  Copyright \copyright{} 2016 \insertshortauthor{}. All Rights Reserved. \texorpdfstring{\\}{}%
  ~ \texorpdfstring{\\}{}%
}
\subject{Intro}

\usepackage{hyperref}

\usepackage{xcolor}
\hypersetup{
	colorlinks,
	linkcolor={red!50!black},
	citecolor={blue!50!black},
	urlcolor={blue!80!black}
}

% bibliography...
\usepackage[backend=biber,bibstyle=numeric-comp,sorting=ydnt,natbib=true,mcite=true,maxnames=100,url=true,isbn=false,doi=false,uniquename=init,giveninits=true,hyperref=true,backref=true,date=edtf,sortcites]{biblatex}
\renewcommand{\subtitlepunct}{\addcolon\addspace}

%\hypersetup{pdfpagemode=FullScreen}


% Load support for hypertext references...
\usepackage{url}

% Turn off the navigation bar...
\beamertemplatenavigationsymbolsempty

% Support arithmetic with numbers...
\usepackage{fp}

% Use Tikz
\usepackage{tikz}

%
% Make it easy to output a number in roman numeral form...
%
\makeatletter
\newcommand*{\asroman}[1]{\expandafter\@slowromancap\romannumeral #1@}
\makeatother

%
% Redefine the part page to subtract one from the part # since the first
% part of this document represents the document title page and overview
% preamble...
%
\setbeamertemplate{part page}
{
  \begin{centering}
    \FPeval{\result}{clip(\insertpartnumber-1)}
    {\usebeamerfont{part name}\usebeamercolor[fg]{part name}\partname~\asroman{\result}}
    \vskip1em\par
    \begin{beamercolorbox}[sep=16pt,center,shadow=false,rounded=true]{part title}
      \usebeamerfont{part title}\insertpart\par
    \end{beamercolorbox}
  \end{centering}
}

%
% Define how beamer shows auto-generated pages...
%
\makeatletter
\newcommand*{\resetsectioncount}{
  \beamer@tocsectionnumber=0\relax%
  \setcounter{section}{0}%
}
\makeatother

\makeatletter
\AtBeginPart{%
  % Restart section numbering within each part...
  \resetsectioncount%
  % And output (if applicable) the part page...
  \ifthenelse{\boolean{showpartpage}}{%
    \frame{\partpage}%
  }{}%
}
\makeatother

\AtBeginSection[]{%
  \ifthenelse{\boolean{showsectiontoc}}{%
    \begin{frame}{Table of Contents}%
      \tableofcontents[sectionstyle=show/hide,subsectionstyle=show/show/hide]%
    \end{frame}%
  }{}%
}

% Bibliography...
\addbibresource{thesis.bib}

% Title Page Logo
\titlegraphic{\includegraphics[width=2cm]{logos/uwin-logo-horz-3col.eps}}

% 
% The presentation slide content follows...
%
\begin{document}
	% Don't show logo or footer on title page...
	{
	\setbeamertemplate{logo}{}
	\setbeamertemplate{footline}{} 
	\begin{frame}[label=TitleSlide,noframenumbering]
	  \titlepage
	\end{frame}
	}
	
	\begin{frame}{Thesis Committee}
		\begin{itemize}
			\item Supervisors: Dr. Richard Frost and Dr. Robert Kent
			\item Internal Reader: Dr. Luis Rueda
			\item External Reader: Dr. Richard Caron
			\item Chair: Dr. Alioune Ngom
		\end{itemize}
	\end{frame}

	\begin{frame}{Outline}
		%\small
		\tableofcontents
	\end{frame}
	
	\sectionwithouttoc{Introduction}

	
	\subsectionwithouttoc{Motivation}
	\begin{frame}{Motivation}
		\begin{itemize}
			\item Internet of Things
			\begin{itemize}
				\item Self driving cars, coffee machines, refrigerators, ovens...
				\item Able to be operated remotely through {\em alternative interfaces}
				\item Accessibility
				\item Natural Language Interfaces
			\end{itemize}
			\item The Semantic Web
		\end{itemize}
	\end{frame}

	\subsectionwithouttoc{The Semantic Web}
	\begin{frame}{The Semantic Web}
		\begin{itemize}
			\item The Semantic Web (part of Web 3.0)
			\begin{itemize}
				\item Distributed and decentralized network of databases
				\item Store facts in the form of {\em triples}
				\item Resource Description Framework (RDF) \cite{w3csemanticweb}
				\item SPARQL Protocol and RDF Query Language (SPARQL)
			\end{itemize}
			\item Can be thought of, informally, as applying the model of the World Wide Web to databases
			\item SPARQL is low-level (like SQL)
		\end{itemize}
	\end{frame}
	
	%\subsectionwithouttoc{Resource Description Framework (RDF)}
	\begin{frame}{Resource Description Framework (RDF)}
		\begin{itemize}
			\item Data is organized into {\em triples}
			\item A {\em triple} is a 3-tuple that has the form \texttt{(subject, predicate, object)}
			\item \texttt{subject}, \texttt{predicate}, and \texttt{object} are Uniform Resource Identifiers (URIs)\cite{w3csemanticweb}.  A URI may denote a {\em name}, {\em location}, or possibly both at once.  Examples:
			\begin{itemize}
				\item URL
				\item ISBN of a book
				\item Telephone number
			\end{itemize}
			\item A database containing these triples is commonly called a {\em triplestore}
			\item Example triple: \texttt{(jane, purchased, pencil\_1)}
			\item Triplestores that describe facts in terms of {\em entities} are called {\em entity-based triplestores}
		\end{itemize}
	\end{frame}
	
	%\subsectionwithouttoc{SPARQL}
	\defverbatim[colored]\lstI{
		\begin{lstlisting}[basicstyle=\ttfamily,keywordstyle=\color{red}]
	PREFIX foaf:  <http://xmlns.com/foaf/0.1/>
	SELECT ?name ?email
	FROM <http://www.w3.org/People/Berners-Lee/card>
	WHERE {
	{      SELECT DISTINCT ?person ?name WHERE { 
	?person foaf:name ?name 
	} ORDER BY ?name LIMIT 10 OFFSET 10    }
	}
		\end{lstlisting}
	}
	
	\begin{frame}{SPARQL Protocol and RDF Query Language}
		\begin{itemize}
			\item RDF Query Language
			\item Similar to SQL
			\item Queries are submitted to {\em SPARQL endpoints}
			\item Form patterns that define restrictions on the elements of triples
			\item Similar to pattern matching in functional programming
			\item Low-level
			\item Example: \lstI
		\end{itemize}
	\end{frame}
	
	\subsectionwithouttoc{The Problem}
	\begin{frame}[allowframebreaks]{The Problem}
		\begin{itemize}
			\item SPARQL is low level and not user-friendly
			\item Natural Language Interfaces
				\begin{itemize}
					\item Require little technical knowledge to use
					\item Require minimal effort on part of the user
					\item Improve accessibility
					\item Area of active research, with several previous attempts (non-trivial)%TODO: HERE
				\end{itemize}
		\end{itemize}
	\end{frame}

	\subsectionwithouttoc{Previous Work}
	\begin{frame}[allowframebreaks]{Previous Work}
		\begin{itemize}
			\item ORAKEL (2007) \cite{cimiano2007orakel}
			\begin{itemize}
				\item Ontology-aware
				\item English interface
				\item Based on Montague semantics
				\item Supports quantification, negation, conjunction
				\item Attempts to directly convert input query into a SPARQL query
			\end{itemize}
			\item QuestIO (2008) \cite{tablan2008natural}
			\begin{itemize}
				\item Ontology-aware
				\item English interface
				\item Keyword oriented
				\item Attempts to match keywords against concepts to hone down queries
				\item Uses SeRQL, similar to SPARQL, to form queries
				\item Sentences transformed into SeRQL queries using formal semantics
			\end{itemize}
		
			\pagebreak
			\item AutoSPARQL (2011) \cite{lehmann2011autosparql}
			\begin{itemize}
				\item Supervised machine learning
				\item Constructs query trees using keywords
				\item Query trees are converted into SPARQL queries
				\item Feedback oriented
				\begin{itemize}
				\item User selects candidates from returned query set that best matches what the user is looking for
				\item User is expected to be actively involved in refining subsequent results
				\end{itemize}
			\end{itemize}
		\end{itemize}
	\end{frame}
	
	\subsectionwithouttoc{Shortcomings of Previous Work}
	\begin{frame}{Shortcomings of Previous Work}
		
		%\textbf{The Problem}:
		Previous work has seen success, but the English NLIs have a shortcoming in that they do not allow for prepositional phrases in queries, and hence have a limited coverage of the English language.
		%\newline
		
	\end{frame}

	\subsectionwithouttoc{New Approach}
	\begin{frame}{New Approach}
		
		The work presented in this thesis draws on two main concepts:
		\begin{itemize}
			\item Executable Attribute Grammars\cite{frosthafiz2008}
			\item Event-Based Denotational Semantics\cite{frostagboola2014}.
		\end{itemize}
	\end{frame}
	
	\subsubsectionwithouttoc{Executable Attribute Grammars}
	\begin{frame}{Executable Attribute Grammars}
		\begin{itemize}
		\item A natural way to process Natural Language queries \cite{frosthafiz2008}
		\item Allow top-down rather than bottom-up parsing, so are highly modular \cite{frosthafiz2008}
		\item Efficient top-down parsing has been demonstrated previously \cite{frosthafiz2008}
		\item In the parser we used:
		\begin{itemize}
			\item Left-recursive grammars are supported directly \cite{frosthafiz2008}
			\item Ambiguity in grammars is tracked, producing parse trees for all possible interpretations of a grammar \cite{frosthafiz2008}
		\end{itemize}
		%\item This allows ambiguous grammars to be used, simplifying development
		\item We modified the parser described in \cite{frosthafiz2008} to support using non-referentially transparent functions as attributes using monads
		\end{itemize}
	\end{frame}
	
	\begin{frame}[allowframebreaks]{Denotational Semantics}
		Introduction to denotational semantics:
		%TODO: Explain quantification see frost2013event prepositions
		\begin{itemize}
			\item Montague Semantics
			\begin{itemize}
				\item Described by Dr. Richard Montague in 1970\cite{dowty2012introduction}
				\item Uses {\em characteristic functions} in {\em higher order logic} as denotations for English words %TODO:
				\item Not computationally tractable due to universal quantification
				\item Examples: %TODO
			\end{itemize}
			\item FLMS
			\begin{itemize}
				\item Described by Frost et al. in 1989\cite{frost1989constructing}, based on Montague's approach
				\item Uses set-theory instead of higher order logic to describe denotations for English words
				\item Includes a denotation for transitive verbs, missing from Montague's approach\cite{frost2013event}
				\item Computationally tractable\cite{frost1989constructing}
				\item Examples: %TODO
			\end{itemize}
		\end{itemize}
	\end{frame}
	
	%\subsubsectionwithouttoc{Event-Based Denotational Semantics}
	\begin{frame}[allowframebreaks]{FLMS on the Semantic Web}
		\begin{itemize}
			\item FLMS can be directly implemented on entity-based triplestores
			\item No prepositions: how to get contextual information from a triple?
			\begin{itemize}
				\item Use quads instead of triples
				\item Reification
				\item {\em Event-Based Triplestores}
			\end{itemize}
		\end{itemize}	
	\end{frame}

	%TODO: Font
	\defverbatim[colored]\lstIoo{
	\begin{lstlisting}[basicstyle=\ttfamily,keywordstyle=\color{red}]
		(jane, purchased, pencil_1)
	\end{lstlisting}
	}

	\defverbatim[colored]\lstIo{
		\begin{lstlisting}[basicstyle=\ttfamily,keywordstyle=\color{red}]
		(event1, subject, Jane)
		(event1, type, purchase)
		(event1, object, pencil)
		\end{lstlisting}
	}

	%\sectionwithouttoc{Event-Based Triplestores}
	\begin{frame}{Event-Based Triplestores}
		\begin{itemize}
			\item Triples describe facts about events
			\item Information about entities and their relationships to one another may be gleaned from the events in which they occur
			\item Additional information about an event may be added by simply adding new triples
			\item Key motivation: directly supports reification on triples\cite{frostagboola2014}
		\end{itemize}
		Example (traditional):
		{\lstIoo}
		Example (event-based):
		{\lstIo}
	\end{frame}

	\subsubsectionwithouttoc{Event-Based Denotational Semantics}
	\begin{frame}[allowframebreaks]{Event-Based Denotational Semantics}
		\begin{itemize}
			\item EV-FLMS 
			\begin{itemize}
				\item Described by Frost et al. in 2013\cite{frost2013event}, based on FLMS
				\item Intended for use with {\em event-based triplestores}
				\item Includes a denotation for prepositional phrases
				\item Key difference with existing approaches: direct translation to SPARQL is not performed
				\item Examples: %TODO
			\end{itemize}
			\item {\em Unified EV-FLMS (UEV-FLMS)}
			\begin{itemize}
				\item Based on EV-FLMS
				\item Solves two problems with how prepositions are treated in EV-FLMS
				\item Novel method of handling the word ``by'', treating it as a preposition
				\item Unifies semantic concepts
				\item Improves efficiency
			\end{itemize}
		\end{itemize}
	\end{frame}
	
	\subsectionwithouttoc{Thesis Statement}
	\begin{frame}{Thesis Statement}
		By integrating a novel event-based denotational semantics with a parser constructed as an executable attribute grammar, it is possible to create a highly modular and extensible Natural Language Interface to the Semantic Web that supports the use of prepositional phrases in queries.
	\end{frame}
	
	\subsectionwithouttoc{Proof of Concept}
	\begin{frame}{Proof of Concept}
		We prove the Thesis by creating an online English query interface to a triplestore containing thousands of facts about the solar system\cite{Solarman:2016}.
		
		Some example queries that can be handled by this system include:
		
		\begin{itemize}
			\item ``when was something discovered at mt\_wilson''
			\item ``how was the thing that was discovered at flagstaff discovered''
			\item ``what was discovered in 1877 in us\_naval\_observatory''
			\item ``what planet is orbited by a moon that was discovered in 1684''
			\item ``which vacuumous moon that orbits jupiter was discovered by nicholson or hall with a telescope in 1938 in mt\_wilson or mt\_hopkins''
		\end{itemize}
		
		The interface can be accessed at: \url{http://speechweb2.cs.uwindsor.ca/solarman2/}
	\end{frame}
	
	\sectionwithouttoc{Demonstration}
	\begin{frame}{Demonstration}
		\begin{itemize}
			\item Demo
		\end{itemize}
	\end{frame}

	%\sectionwithouttoc{Implementation Details}
	\begin{frame}{Implementation Details}
		\begin{itemize}
			\item Implemented in Haskell
			\item SafeHaskell is used in the Direct Query Interface to prevent arbitrary code execution
			\item The monadic parser was joined with the monadic semantics in combination with a grammar
			\item Memoization is performed at the query level to keep queries fast
			\item The actual query logic is contained in the Getts module, which provides interfaces to SPARQL triplestores and also list-based in-program triplestores.  A general typeclass is provided for future extensibility
		\end{itemize}
	\end{frame}

	\sectionwithouttoc{Unified EV-FLMS}
	\begin{frame}[allowframebreaks]{Unified EV-FLMS}
	\end{frame}
	
	%\sectionwithouttoc{Improvements over Original Semantics}
	\begin{frame}{Improvements over Original Semantics}
		\begin{itemize}
			\item The original event-based semantics were developed by Frost, Agboola et al. \cite{frostagbola2014}
			\item The interface to the remote triplestore was overhauled with the monadic improvements described
			\item The original semantics suffered from an implicit `and' problem where chained prepositions in sentences were treated as though they had the word `and' in between them.
			For example ``who discovered a moon in 1877 with a telescope'' would have been treated as ``who discovered a moon in 1877 and with a telescope''.
			\item This also interfered with the correct handling of words such as `one', `two', and `every'
			\item The `collect' operation is now more efficient an can be performed in O(n log n) time
			\item `by' is treated as a preposition in the grammar
		\end{itemize}
	\end{frame}
	
	%\sectionwithouttoc{Novel additions}
	\begin{frame}{Novel additions}
		The use of `by' as a preposition:
		\begin{itemize}
		\item In developing the system, it was discovered that the logic for handling the word `by' (e.g, ``phobos was discovered by hall'') was identical to a preposition in the semantics that searched for a subject in a given event
		\item The functionality was merged, and with no added complexity to the grammar, the word `by' is supported in queries
		\end{itemize}
	\end{frame}
		
	\sectionwithouttoc{Timing}
	\begin{frame}{Timing}
		\begin{itemize}
			\item Two experiments
			\item Results:
		\end{itemize}
	\end{frame}

	\sectionwithouttoc{Conclusions and Future Work}
	\begin{frame}{Conclusions and Future Work}
		\begin{itemize}
			\item IO is clearly the limiting factor: many small queries produced, each query has overhead
			\item Scaling up:
			\begin{itemize}
				\item More advanced Query Fusion using monads
				\item Exploiting data-parallelism
				\item Foregoing set-theory and delving into conceptual spaces
			\end{itemize}
			\item Compatibility:
			\begin{itemize}
				\item Providing event based views of entity-based triplestores exploiting ontology data
			\end{itemize}
		\end{itemize}
	\end{frame}
	
	\sectionwithouttoc{Thank You \& Feedback}
	\begin{frame}{Thank You \& Feedback}
		\centering
		Thank you for attending! \\ Comments, suggestions, or questions?
	\end{frame}
  
	\sectionwithouttoc{References}
	\begin{frame}[allowframebreaks]{References}
	  \def\newblock{}
	  \setbeamertemplate{bibliography item}[text]
	  \setlength\bibitemsep{1.5\itemsep}
	  \printbibliography
	\end{frame}
	
	%\appendix 
	%\section*{Appendix Overview}
	%\begin{frame}{Appendix Overview}
	%  \tableofcontents[part=3,sectionstyle=show/show,subsectionstyle=hide/hide/hide]
	%\end{frame}
	
	%\sectionwithouttoc{An Appendix Title}
	%\begin{frame}{\insertsection}
	%  This is the first section.
	%\end{frame}
	
	%\sectionwithouttoc{Another Appendix Title}
	%\begin{frame}{\insertsection}
	%  This is the second section.
	%\end{frame}
\end{document}
